% Bericht -- Gerüst der Harwness-Vorlage (Skill latex-report).
% Bauen: latex.build mit engine = xelatex; latex.build läuft so oft wie nötig
% (Inhaltsverzeichnis). Alle Texte unten sind kurze Platzhalter: ersetzen,
% kürzen oder streichen. Nichts erfinden -- was der Auftrag nicht hergibt,
% bleibt als % TODO: stehen und wird in der Übergabe genannt.
\documentclass[11pt,a4paper]{scrartcl}

% ── Anpassung: von latex.template gefüllt, Vorgaben in harw-report.sty ──
\newcommand{\harwLanguage}{%%LANGUAGE%%}
\newcommand{\harwAccentHex}{%%ACCENT%%}
\newcommand{\harwWarnHex}{%%WARN%%}
\newcommand{\harwMainFont}{%%MAINFONT%%}
\newcommand{\harwSansFont}{%%SANSFONT%%}
\newcommand{\harwMonoFont}{%%MONOFONT%%}
\usepackage{harw-report}

\title{%%TITLE%%}
\subtitle{%%SUBTITLE%%}
\author{%%AUTHOR%%}
\date{%%DATE%%}

\begin{document}
\maketitle

\begin{achtung}[Wichtig vorab]
Das Wichtigste, was die Leserin vor allem anderen wissen muss, in zwei bis
drei Sätzen: eine Einschränkung, ein Status oder eine offene Entscheidung.
Abschnitt~\ref{sec:stand} beschreibt den Stand genauer.
\end{achtung}

\begin{merke}[Die Idee in drei Sätzen]
Erster Satz: worum es geht. Zweiter Satz: was dabei herauskommt. Dritter
Satz: was die Leserin daraus machen kann.
\end{merke}

\tableofcontents
\newpage

\section{Worum es geht}\label{sec:ueberblick}

Ein kurzer Absatz zu Anlass, Ziel und Leserkreis des Berichts. Dateien und
Pfade stehen in \code{\textbackslash path}, etwa \path{daten/messwerte.csv};
kurze Bezeichner in \code{\textbackslash code}, etwa \code{status}.

\subsection{Kleines Glossar}

\begin{tabularx}{\textwidth}{@{}lL@{}}
\toprule
\textbf{Begriff} & \textbf{Bedeutung} \\
\midrule
Begriff A & Eine kurze Erklärung in einem Satz. Längere Texte umbrechen in
            der L-Spalte von selbst. \\
Begriff B & Eine zweite Erklärung. \\
\bottomrule
\end{tabularx}

\section{Aufbau}\label{sec:aufbau}

Ein Satz, der die Abbildung einordnet: Was soll die Leserin darin sehen?

\begin{figure}[htbp]
\centering
\begin{tikzpicture}[node distance=1.6cm]
  \node[harwbox,fill=harwgrau] (eingang) {Eingang\\(Daten)};
  \node[harwbox,fill=harwaccent!15,right=of eingang] (kern)
    {\textbf{Verarbeitung}\\prüfen, aufbereiten};
  \node[harwbox,fill=harwgrau,right=of kern] (ergebnis) {Ergebnis\\(Bericht)};
  \node[harwbox,fill=harwgrau,below=1.2cm of kern] (ablage) {Ablage\\(Archiv)};
  \draw[harwpfeil] (eingang) -- node[harwlabel,above=2pt]{liefert} (kern);
  \draw[harwpfeil] (kern) -- node[harwlabel,above=2pt]{erzeugt} (ergebnis);
  \draw[harwdoppelpfeil] (kern) -- node[harwlabel,right=2pt]{speichert,\\liest} (ablage);
\end{tikzpicture}
\caption{Die Bausteine und ihre Verbindungen (Beispiel)}\label{fig:aufbau}
\end{figure}

\section{Ergebnisse im Detail}\label{sec:ergebnisse}

Ein Satz zur Tabelle. Lange Tabellen (ab etwa 15 Zeilen) als
\code{longtable}; der Kopf wiederholt sich auf jeder Seite.

{\small
\begin{longtable}{@{}>{\raggedright}p{3.0cm}>{\raggedright}p{4.4cm}>{\raggedright\arraybackslash}p{7.6cm}@{}}
\toprule
\textbf{Bereich} & \textbf{Befund} & \textbf{Bedeutung} \\
\midrule
\endhead
\bottomrule
\endlastfoot
Bereich 1 & Kurzer Befund. & Was der Befund für die Leserin heißt. \\
Bereich 2 & Kurzer Befund. & Was der Befund für die Leserin heißt. \\
Bereich 3 & Kurzer Befund. & Was der Befund für die Leserin heißt. \\
\end{longtable}
}

\begin{beispiel}
Ein kurzes, konkretes Beispiel macht einen abstrakten Befund greifbar.
\end{beispiel}

\section{Stand und nächste Schritte}\label{sec:stand}

\begin{description}[leftmargin=2.6cm,style=nextline,font=\sffamily\bfseries\small]
  \item[Stand] Was fertig ist und belegt werden kann.
  \item[Offen] Was noch fehlt oder ungeklärt ist.
  \item[Nächster Schritt] Wer was bis wann tut.
\end{description}

\begin{thebibliography}{9}
\addcontentsline{toc}{section}{\refname}
\small
\bibitem{quelle1} Nachname, Vorname: \emph{Titel der Quelle}. Verlag, Jahr.
  \url{https://example.org/quelle}
\end{thebibliography}

\end{document}
